\section{All linear Weil representations in the transitive normal-core gate}
\label{sec:linear-weil-core}

\paragraph{Scope and source distinction.}
Here a linear Weil representation means a nontrivial constituent of the
complex permutation representation of $\operatorname{GL}_n(q)$ on
$\mathbf F_q^n$, restricted to $\operatorname{SL}_n(q)$, with $n\ge3$.
This is the general-linear family in G\'erardin's Proposition~4.2,
pp.~82--83, including its complete function-space proof
\cite{BI:GerardinLinear}. It is different from the symplectic and
unitary Weil families. The model is published; the support estimates,
source exclusion and finite endpoint computation below are new
derivations. No non-Weil representation of a linear group is excluded
merely by this section.

As before, $J$ is the finite image of an inducing representation
$\tau$ for the six-dimensional multiplicity variation, and
$e=\dim_{\mathbf C}\tau$. Its character field may be larger than the
original imaginary quadratic field $K$; only
$K\subset\mathbf Q(\chi_\tau)$ is used. Write
$\nu(T)=\min_z\operatorname{rank}(T-z)$.

\begin{lemma}[D-BI-23a: models, irreducibility and trace bounds]
\label{D-BI-23a}
Put
\[
 D=\frac{q^n-1}{q-1},\qquad
 A=D-1=\frac{q^n-q}{q-1},\qquad
 B=\frac{q^{n-1}+q-2}{q-1}.
                                                               \tag{LW.1}
\]
For every nontrivial character $\vartheta$ of $\mathbf F_q^*$,
the space
\[
 W_\vartheta=\{f:\mathbf F_q^n\setminus\{0\}\longrightarrow\mathbf C:
                f(tx)=\vartheta(t)f(x)\}
                                                               \tag{LW.2}
\]
is irreducible under $\operatorname{SL}_n(q)$ and has degree $D$.
The deleted projective-point permutation space $W_1^0$ is irreducible
of degree $A$. For a nonscalar element of the represented core,
\[
 |\chi_\vartheta(g)|\le B\quad(\vartheta\ne1),\qquad
 |\chi_1^0(g)|\le B-1.                                       \tag{LW.3}
\]
When $q=2$ the second bound improves to
$|\chi_1^0(g)|\le2^{n-1}-2$.
Every model with $q=2$ or $q=3$ is realizable over $\mathbf R$;
$W_1^0$ is real for every $q$.
\end{lemma}

\begin{proof}
G\'erardin's model and degree calculation give (LW.2) and the deleted
projective permutation summand. We check restriction to the special
linear group, rather than assuming it preserves irreducibility.
For $n\ge3$, $\operatorname{SL}_n(q)$ is transitive on ordered
linearly independent pairs: extend the pair to a basis and adjust
the third basis vector to correct the determinant. Its orbits on
$\mathbf F_q^n\times\mathbf F_q^n$ are consequently the independent
pairs, the origin, the two types with exactly one vector zero, and
the $q-1$ types $(x,tx)$ with $x\ne0$. There are $q+3$ orbits.

The permutation representation on $\mathbf F_q^n$ decomposes as
two trivial summands, $W_1^0$, and the $q-2$ nontrivial spaces
$W_\vartheta$. Each of the latter spaces is nonzero. The squared
character norm is exactly the number $q+3$ of orbits just computed,
which is already the minimum possible contribution
$2^2+1+(q-2)$. Thus all the indicated nontrivial summands are
irreducible and pairwise inequivalent upon restriction. This also
handles $q=2$, when there are no nontrivial $\vartheta$.

For $\alpha\in\mathbf F_q^*$ let
$d_\alpha=\dim\ker(g-\alpha)$. The scalar-character projector on
the nonzero-vector permutation model gives, up to replacing
$\vartheta$ by its inverse,
\[
 \chi_\vartheta(g)=\frac1{q-1}
        \sum_{\alpha\in\mathbf F_q^*}
                  \vartheta(\alpha)(q^{d_\alpha}-1).
                                                               \tag{LW.4}
\]
For $\vartheta=1$, the same sum counts fixed projective points;
one subtracts one for the deleted summand.

If $g$ is not a scalar natural matrix, then
$d_\alpha\le n-1$ and $\sum d_\alpha\le n$.
For nonnegative integers with these constraints, convexity of
$q^d-1$ bounds their sum by
$q^{n-1}-1+(q-1)=q^{n-1}+q-2$.
For example, transferring a unit from a smaller positive part to
a larger part weakly increases the sum until the pattern is
$(n-1,1)$, or a smaller total pattern. This proves the first
bound. The deleted projective character is an integer between
$-1$ and $B-1$; here $B-1\ge1$, proving the second.
When $q=2$ there is only one possible nonzero scalar eigenvalue,
so its exact character is $2^{d_1}-2$. Since $d_1\le n-1$ and
$n\ge3$, its absolute value is at most $2^{n-1}-2$.

These estimates apply to every nonscalar element of the represented
core. Indeed scalar natural matrices act scalarly on each model.
Conversely the preimage of the projective kernel is a normal
subgroup of $\operatorname{SL}_n(q)$. Its image in the simple group
$\operatorname{PSL}_n(q)$ is trivial or the whole group. The latter
would make a nontrivial irreducible representation of degree
greater than one projectively scalar, which is impossible.
Thus the projective kernel is exactly the natural center. This
argument also shows that any central quotient afforded by the
model has the same relevant projective core.

The projective permutation model is real. For $q=3$ the only
nontrivial scalar character has values $\pm1$, and (LW.2) has
a basis of real-valued functions. These give the asserted real
forms, and averaging a real inner product supplies a nondegenerate
invariant symmetric bilinear form after complexification.
\end{proof}

\begin{lemma}[D-BI-23b: the degree-twenty-one endpoint]
\label{D-BI-23b}
In either degree-twenty-one linear Weil representation of
$U=\operatorname{SL}_3(4)=3.\operatorname{L}_3(4)$, every
nonscalar finite-order linear normalizer element satisfies
\[
                         \nu(T)\ge8.                         \tag{LW.5}
\]
All scalar phases and all compatible outer automorphisms are included.
\end{lemma}

\begin{proof}
The two models have faithful central characters of the order-three
center, and hence are faithful on $U$. The published ATLAS data give
\[
 |\operatorname{L}_3(4)|=20160,\quad
 \operatorname{Mult}(\operatorname{L}_3(4))=C_3\times C_4\times C_4,
 \quad\operatorname{Out}(\operatorname{L}_3(4))=D_{12}
          \simeq C_2\times S_3.                              \tag{LW.6}
\]
The primary ATLAS page supplies the natural three-dimensional
matrices of the triple cover \cite{BI:AtlasL34}, but also carries
a general work-in-progress warning. We therefore check the
structural assertions independently. Counting ordered bases gives
$|\operatorname{GL}_3(4)|=63\cdot60\cdot48=181440$; the
determinant quotient and special-linear center both have order
three, giving $|U|=60480$ and the displayed simple-group order.
The exceptional status is confirmed in \cite[Chapter~5,
Table~1, p.~35]{BI:HHM}; that table lists the exceptional group,
rather than purporting to specify the multiplier structure.
The full structure $C_4^2\times C_3$ is independently specified
and used in Breuer's complete Schur-cover construction,
\S2.7-4, explicitly citing the printed ATLAS p.~23
\cite{BI:BreuerL34}.

Breuer's \S2.6-4 independently records the full outer group
$D_{12}$ and its three order-two subgroup types. This is also
consistent with the standard special-linear automorphism
description in his \S11.5-7: the diagonal outer generator
$\delta$ has order three, the field involution $f$ squares
matrix entries, and the graph involution $\gamma$ is inverse
transpose. The last two commute and both invert $\delta$.
Thus $(C_3\rtimes C_2)\times C_2\simeq S_3\times C_2$
is the stated outer group. The numerical character-table
premises below are the complete printed-ATLAS library records,
not an inference from the warning-bearing web page.
In particular $U$ is the actual triple cover,
not an assumed universal cover. The prime-to-three part of the
universal center is characteristic, so every simple-group
automorphism lifts to $U$. An automorphism of the perfect group
$U$ trivial on its central quotient is trivial: its central
discrepancy is a homomorphism to an abelian group.

Every cyclic subgroup of $D_{12}$ is conjugate to the identity,
one of its three order-two types, its order-three subgroup, or its
order-six subgroup. We retain the complete six published ATLAS
tables and the five actual normal-core fusions:
\[
\begin{array}{c|r|l|r|r}
\text{table}&\text{classes}&\text{selected rows}&
       \text{nonscalar entries}&\min\nu\\ \hline
3.L_3(4)&28&17,18&50&8\\
3.L_3(4).2_1&38&27,28,29,30&140&8\\
3.L_3(4).2_2&23&\text{none}&0&\text{--}\\
3.L_3(4).2_3&23&\text{none}&0&\text{--}\\
3.L_3(4).3&60&25,\ldots,30&342&8\\
3.L_3(4).6&66&31,\ldots,42&756&8.
\end{array}                                                   \tag{LW.7}
\]
A row is selected precisely when its pointwise restriction along
the retained fusion equals one of the two degree-twenty-one core
rows. Selection is not based only on the degree. Each table record
explicitly credits the ATLAS as its origin. The input file
\texttt{linear\_weil/linear\_weil\_inputs.json} contains every
character row, class size and order, all prime power maps, the
five fusions, and the source metadata. Its SHA-256 is
\texttt{12becae3b988a84d6820b658da1ac3dcba4a2a83ff3be85c1342e6cd9bed5e76}.
This is copied published numerical input; the following reconstruction
is an exact independent arithmetic implementation of its consequences.

For a class of order $d$, its eigenvalue multiplicities are
\[
 m_k=\frac1d\sum_{r=0}^{d-1}\chi(g^r)\zeta_d^{-kr},
                  \qquad0\le k<d.                           \tag{LW.8}
\]
The prime power maps and cyclotomic Galois action determine every
power in this expression. The GAP program checks all table-row
orthogonality and computes all spectra. The separate Python program
\texttt{check\_linear\_weil.py} selects rows by polynomial
restriction equality, verifies all five normal fusions and cyclic
quotients, and reconstructs (LW.8) by normalized rational traces,
that is, Ramanujan sums. It imports no GAP scan results. Both
programs pass all twenty-four selected rows and all1288 nonscalar
entries. The minimum moving rank is eight in each selected row.

We explain why this finite calculation covers an arbitrary actual
normalizing operator; no assertion that the whole normalizer equals
one tabulated extension is necessary. Let $\alpha$ be the induced
automorphism of $U$. Its outer class generates one of the cyclic
subgroups just listed. After conjugating that subgroup and changing
$\alpha$ by an inner automorphism, use a corresponding extension
$A=\langle U,t\rangle$ from (LW.7), with $A/U$ cyclic of order $d$.
The transformed core representation is still one of the two
degree-twenty-one rows. The actual operator implementing the
automorphism differs from the chosen intertwiner only by a scalar.
The relation
\[
            T^d\rho(t^d)^{-1}\in\mathbf C^*
                                                               \tag{LW.9}
\]
follows from Schur's lemma. Rescale $T$ by a $d$th root so that the
defining relation holds exactly. Then $\rho$ extends to an actual
representation of $A$, irreducible because its restriction to $U$
is irreducible. It therefore occurs in one of the selected rows.
The original element, including its possible inner factor, is scalar
times a matrix of some element of $A$, all of whose classes have
been checked. Scalar rescaling preserves $\nu$. If the relevant
table has no selected row, no such compatible operator exists.
The case of trivial outer class is covered by the core table.
This proves (LW.5).
\end{proof}

\begin{theorem}[D-BI-23: exclusion of transitive linear Weil cores]
\label{D-BI-23}
Suppose the inducing representation of an actual finite-cover
Weil-sixfold source has an irreducibly acting normal central product
of quasisimple factors whose simple quotients are
$\operatorname{PSL}_n(q)$, with $n\ge3$. Suppose every factor
representation is a linear Weil model of D-BI-23a, and the full
inducing group permutes the factors transitively. Then its period
map has rank less than nine.

This includes the transitive forced-core situation of D-BI-19,
all compatible outer automorphisms, arbitrary scalar phases and
tensor-factor permutations. It allows any imaginary quadratic
original field, a larger inducing character field, and either
polarization determinant. It does not exclude non-Weil factor
characters or claim a unitary-family result.
\end{theorem}

\begin{proof}
Transitivity makes the factor degrees equal, say $a$, and the total
degree is $e=a^t$ for some $t\ge1$. Each $a$ is $A$ or $D$, and
$a\ge6$. A nonscalar element of the represented central product
has at least one nonscalar factor, so a factor character bound
$|\chi(g)|\le ca$ implies the same normalized bound $c$ on the
entire core. D-BI-18 therefore applies to every nonscalar element
of the full inducing group, including those permuting factors.

For completeness, the actual intermediate-cover Euler identity is
$6=2e(h-1)+\sum_j\operatorname{rank}(1-g_j)$. If $h\ge2$, its
rightmost sum would be $6+2e(1-h)<0$. If $h=1$, there are at most
six genuine branch points and at most six moduli parameters.
Rank nine consequently requires $h=0$, at least twelve genuine
inertias, and total moving rank $2e+6$. These are consequences of
the actual period factorization, not a claim that the dimension
inequality suffices for dominance.

First let $q\ge5$. In (LW.3) use the common bound
$c=B/A<1$. Since $e\ge A$, the necessary continuous inequality
of D-BI-18 yields
\[
 6\ge e(1-3c)\ge A-3B
       =\frac{(q-3)q^{n-1}-4q+6}{q-1}>6.                   \tag{LW.10}
\]
The last expression increases with $n$ for $q\ge5$. At $n=3$
it equals $q^2-2q-6$, which is nine at $q=5$ and strictly
increases thereafter. This is a contradiction.
For $q=4$ and $n\ge4$, the same expression is
$(4^{n-1}-10)/3$, at least eighteen, also a contradiction.

It remains to consider $q=2,3$, and $(q,n)=(4,3)$.
In every real factor case of D-BI-23a, the full irreducible tensor
core has a unique invariant symmetric bilinear-form line. Its
normalizer is essentially self-dual. The inclusion
$K\subset\mathbf Q(\chi_\tau)$ forces a nontrivial similitude
character: if the form were preserved exactly, the character would
be real. Since genuine inertias generate and have product one,
at least two have nontrivial multiplier, and their fixed spaces
are totally isotropic. D-BI-18a gives the necessary bound
\[
                         e(3-5c)\le12.                     \tag{LW.11}
\]
This use of the original field does not identify it with the
possibly larger inducing field.

For $q=3$ and the nontrivial scalar character, $a=D\ge13$ and
\[
 B=\frac{a+2}{3},\qquad
 a(3-5B/a)=\frac{4a-10}{3}\ge14>12.
\]
Since $e\ge a$ and $3-5B/a>0$, (LW.11) is impossible.
For the deleted summand, $a=A\ge12$ and $B-1=a/3$, so it is
already impossible because $a(3-5/3)=4a/3\ge16$.

For $q=2$, there is only the real deleted summand,
$a=2^n-2$, and the improved bound gives
$c=(a-2)/(2a)$. Hence
\[
                  e\le\frac{24a}{a+10}.                    \tag{LW.12}
\]
For $n\ge5$, $a\ge30$ and even $e=a$ violates (LW.11), since
$a(3-5c)=a/2+5>12$. For $n=3,4$, (LW.12) is respectively
$e\le9$ and $e\le14$, forcing $t=1$ and $a=6,14$.
The first core is $\operatorname{SL}_3(2)\simeq
\operatorname{PSL}_2(7)$; its degree-six full normalizer has
$\nu\ge a/4$ by D-BI-21c. The second is
$\operatorname{SL}_4(2)\simeq A_8$, and its ordinary degree-fourteen
full normalizer has the same bound by HD-ALT-02a(3).
Both are incompatible with the exact branch budget:
\[
           12\lceil6/4\rceil=24>18,
       \qquad12\lceil14/4\rceil=48>34.                      \tag{LW.13}
\]
These endpoint spectral bounds themselves have no nonsplit
polarization hypothesis.

Finally take $(q,n)=(4,3)$. For the deleted summand, $a=20$,
$B-1=5$ and $c=1/4$. It is real, and
$a(3-5c)=35>12$, excluding it and all its transitive tensor powers.
For either nontrivial scalar character, $a=21$, $B=6$ and
$c=2/7$. D-BI-18 forces $e\le42$, hence $t=1$.
D-BI-23b gives actual moving rank at least eight for every
nonscalar inertia; a nonidentity scalar moves all twenty-one
dimensions. Thus total moving rank is at least
$12\cdot8=96$, whereas $2e+6=48$. This last contradiction
completes every prime-power and every tensor-exponent case.
\end{proof}

\paragraph{Reproducibility and limits.}
Run \texttt{linear\_weil/export\_and\_scan.g} in the retained GAP
runtime and \texttt{python check\_linear\_weil.py} in this directory.
The latter also checks the displayed endpoint arithmetic. Its first
run stopped on a metadata-string expectation: the library writes
the multiplier as \texttt{(3x4\^{}2)}, not \texttt{3.4\^{}2}.
Correcting that literal expectation gave a complete PASS; no
mathematical inequality or table entry changed. Source-integrity
checks are recorded separately from TeX compilation, which was
unavailable. The transitive-core hypothesis is explicit and is
supplied in the main primitive tensor-indecomposable gate by
D-BI-19. No arbitrary mixed-type tensor classification is asserted.
